\section{Methodology}
\label{sec:methodology}
To simulate realistic judgment prediction and evaluate the role of RAG in enhancing legal decision-making, we design a modular experimental setup. This setup explores how different types of legal information, such as factual summaries, statutes, and precedents affect model performance on the dual tasks of prediction and explanation. To ensure reproducibility and transparency, we detail the full experimental setup, including model configurations, training routines, and task-specific hyperparameters. The details are mentioned in Appendix~\ref{sec:Experimental-setup}. This includes separate sections for explanation generation (summarization) and legal judgment prediction tasks, outlining all relevant decoding strategies, optimization settings, and dataset splits used across our pipeline variants.

\subsection{Pipeline Construction}
To systematically evaluate the impact of legal knowledge sources, we constructed multiple input pipelines using combinations of the dataset components described in Section~\ref{sec:dataset}. Each pipeline configuration represents a distinct input scenario reflecting different degrees of legal context and retrieval augmentation. These pipelines are as follows:

\begin{itemize}
    \item \textbf{CaseText Only:} Includes only the summarized version of the full case judgment, which contains factual background, arguments, and reasoning.

    \item \textbf{CaseText + Statutes:} Appends summarized statutory references cited in the judgment to the case text, simulating scenarios where relevant laws are explicitly considered.

    \item \textbf{CaseText + Precedents:} Incorporates prior cited judgments mentioned in the original case, representing explicitly relied-upon precedents.

    \item \textbf{CaseText + Previous Similar Cases:} Adds top-3 semantically similar past judgments (retrieved via ChromaDB using \texttt{all-MiniLM-L6-v2} embeddings), allowing the model to learn from precedents not explicitly cited.

    \item \textbf{CaseText + Statutes + Precedents:} A comprehensive legal input pipeline combining the full judgment summary, statutes, and cited prior judgments.

    \item \textbf{Facts Only:} A minimal pipeline containing only the factual summary, excluding all legal reasoning and verdicts. This setup evaluates whether a model can infer judgments from facts alone.

    \item \textbf{Facts + Statutes + Precedents:} Combines factual input with statutory and precedent context to simulate realistic courtroom conditions where judges rely on facts, applicable law, and relevant past cases.
\end{itemize}

This modular design enables granular control over input features and facilitates direct comparison of how each knowledge source contributes to judgment prediction and explanation generation.

\subsection{Prompt Design}
\label{subsec:prompts}
To ensure consistency and interpretability across all pipelines, we used fixed instruction prompts with minor variations depending on the available contextual inputs (e.g., facts only vs. facts + law + precedent). These prompts guide the model in producing both binary predictions and natural language explanations. Prompts were structured to reflect real judicial inquiry formats, aligning with the instruction-following capabilities of modern LLMs. Full prompt templates are listed in Appendix Table~\ref{tab:judgment_prediction_prompts_few}, along with prediction examples.

\subsection{Inference Setup}
\label{subsec:inference}
We use the \texttt{LLaMA-3.1 8B Instruct}~\cite{dubey2024llama} model for all experiments in a few-shot prompt setup. Each input sequence, composed according to one of the pipeline templates, is paired with a relevant prompt that asks to output:

\begin{itemize}
    \item A binary judgment prediction: {0} (appeal rejected) or {1} (appeal fully/partially accepted)
    \item A justification: a coherent explanation based on legal facts, statutes, and precedent
\end{itemize}

The model is explicitly instructed to reason with the provided information and emulate judicial writing. Retrieved knowledge (via RAG) is included in-context to enhance legal reasoning while minimizing hallucinations.

% \noindent
This experimental design allows us to evaluate the effectiveness of legal retrieval and summarization under realistic judicial decision-making constraints in the Indian common law setting.



% This lightweight yet powerful setup ensures model reproducibility, maintains privacy (no training on sensitive data), and supports plug-and-play deployment in retrieval-enhanced legal applications.























% \section{Methodology}
% \subsection{Pipeline Creation}

% To evaluate the impact of different legal information types on model performance, multiple input pipelines were constructed by selectively combining components from the curated dataset. Each pipeline represents a unique configuration of textual content derived from Indian legal documents. The base dataset was first segmented into modular components, namely: \textit{CaseText}, \textit{Cited Cases}, \textit{Statutes}, \textit{Previous Similar Cases}, and \textit{Facts}.

% The pipelines were constructed as follows:

% \begin{itemize}
%     \item \textbf{CaseText\_only}: Includes the full judgment text of a case. This forms the base representation encompassing factual background, legal arguments, citations, and judicial reasoning in a summarized format.

%     \item \textbf{CaseText + Statutes}: Combines the summarized case text with the summarized statutory provisions referenced in the document. Statutes were extracted from structured metadata (e.g., “Cited Laws”) and appended in a standardized format.

%     \item \textbf{CaseText + Cited Cases}: Merges the summarized judgment with its explicitly cited prior cases. Cited judgments were identified from the document body, and structured references were scraped from the original source.

%     \item \textbf{CaseText + Previous Similar Cases}: Augments the summarized judgment with the top three semantically similar cases, identified using vector-based similarity search across the entire corpus. This semantic similarity was computed using sentence embeddings and managed through ChromaDB.

%     \item \textbf{CaseText + Statutes + Cited Cases}: An enriched pipeline that combines the summarized text of the case with both statutory references and cited judgments, providing a comprehensive legal context.

%     \item \textbf{Facts Only}: Contains only the factual portion of each case. This subset excludes legal reasoning, verdicts, or citations and serves as a baseline focused solely on the narrative of the case.

%     \item \textbf{Facts + Statutes + Cited Cases}: Integrates factual content with relevant statutes and cited cases, isolating the effect of non-reasoning-based legal inputs on downstream tasks.

% \end{itemize}

% This design ensures comparability across experiments and facilitates systematic evaluation of the influence of different legal content types on the model's performance.


% \subsection{Prompts Used}
% To ensure a consistent and interpretable model behavior across all pipeline configurations, a fixed prompt was used during inference with slight modification based on the supporting information along with case judgment. The prompt acts as an instruction or query to guide the model's response based on the input content. This approach aligns with the prompt-based learning paradigm used in instruction-tuned language models. The details of these prompts can be found in Table \ref{tab:judgment_prediction_prompts_few} in the Appendix. The relevant details and examples of the generated predictions and explanations are also available in the referenced tables in the Appendix.

% \subsection{Inference Setup}

% All model inferences were conducted using the \texttt{Llama 3.1} model in a few-shot setting. For each input, the model received one of the seven instruction templates and was prompted to return:
% \begin{itemize}
%     \item A binary judgment decision (\texttt{1} or \texttt{0})
%     \item A natural language explanation supporting that decision
% \end{itemize}

% The model was explicitly instructed to reason from the legal information provided and justify its conclusion in a format resembling legal rationale.

